\chapter{Conclusion \& Outlook} \label{ch:conclusion}
This thesis assessed whether the clustering of FRBs can constrain the bias of their
host population, and to what extent cross-correlating them with a photometric galaxy survey can improve this measurement.
A complete forecasting pipeline was developed that combines the nonlinear matter power spectrum with radial 
weight functions for modeled FRB samples and the six tomographic bins of the KiDS-Legacy data release to 
obtain angular auto- and cross-power spectra using the Limber approximation.
Poissonian shot noise was added to the auto-spectra of both tracer populations, 
while the cross-spectrum was assumed to be free of shot noise. The resulting spectra were then 
propagated into a Fisher forecast for the FRB host-bias amplitude $b_0$ and evolution exponent $\delta$.

The combination of the FRB sample with the KiDS galaxy survey yields strong constraints on the 
parameters of interest. The marginalized uncertainties are reduced by factors of $17$ to $51$, while 
the FoM improves by factors of $61$ to $449$, despite using a footprint roughly $28$ 
times smaller than that assumed for the FRB-only analysis. These results demonstrate the strong 
constraining power gained from combining the FRB and galaxy samples. However, the current KiDS 
footprint is not yet sufficient to distinguish between a prompt magnetar channel and a delayed 
neutron-star channel, as the separation between the two fiducial models remains well within 
the $1\sigma$ contours.

\section{Limitations} \label{sec:limitations}
The analysis rests on some simplifying assumptions that may affect the forecast.

\begin{itemize}
    \item The Limber approximation restricts the analysis to multipoles $\ell \geq 10$,
    excluding precisely the large angular scales.
    \item The Gaussian likelihood and purely Gaussian covariance neglect non-Gaussian contributions
    from nonlinear structure growth. The forecast uncertainties should therefore be interpreted as
    optimistic lower bounds rather than as realistic estimates.
    \item The FRB redshift distribution and fiducial bias values are model assumptions rather than
    measured quantities.
    \item The survey geometry is simplified by accounting for an incomplete sky coverage only
    through the $f_{\mathrm{sky}}$ prefactor, while a real analysis would require a full
    pseudo-spectrum treatment~\cite{Hivon_2002}.
\end{itemize}

We use a power law, $b = b_0(1+z)^{\delta}$, which provides a reasonable 
description of the locally observed FRB host population at low redshift~\cite{Berti_2021}. 
At higher redshift, the bias of galaxy populations is expected to evolve more strongly and to 
depend on the host halo mass~\cite{Pan_2020,Tinker_2010}, making a simple power-law description 
less accurate. A description valid over the full redshift range would require a halo occupation 
distribution (HOD) model, which is beyond the scope of this work. For our purposes, however, the power 
law provides an adequate approximation, as the constraints are dominated by $z \lesssim 1.5$.

DAS SCHOENER EINPFLEGEN

\section{Outlook} \label{sec:outlook}
Since the Fisher information scales linearly with $f_{\mathrm{sky}}$, the most direct improvement
is to replace the KiDS field with a wider galaxy survey. The Legacy Survey of Space and Time
(LSST) at the Vera C. Rubin Observatory will image about $\qty{18000}{\deg\squared}$ of the
southern sky~\cite{Ivezic_2019}, corresponding to $f_{\mathrm{sky}} = 0.4363$, roughly 13
times the KiDS area. Running the pipeline with the ten-bin Y10 lens sample of the LSST Dark Energy Science Collaboration
shown in \autoref{fig:lsst_y10_nz}~\cite{Zhang_2025_LSST}
while keeping all other ingredients unchanged yields the contours shown in
\autoref{fig:fisher_kids_vs_lsst} and the numbers listed in \autoref{tab:lsst_sigmas} and \autoref{tab:lsst_fom}.

The improvement is substantial: for the deep magnetar configuration, the uncertainty on $b_0$ drops from
$1.3$ to $0.42$. The FoM increases by factors of $9.1$--$12.0$, slightly less than the factor of $13$
expected from the larger footprint alone, since the FoM scales linearly with
$f_{\mathrm{sky}}$. Per unit area, the LSST lens sample is thus slightly less informative than KiDS,
most notably for the deep survey, because the LSST lens bins end at $z \approx 1.2$, whereas the sixth
KiDS bin extends beyond $z = 1.5$ into the survey's high-redshift tail.
The separation between the two progenitor models, $\Delta b_0 = 0.5$ and
$\Delta\delta = 0.6$, becomes comparable to the marginalized uncertainties in the deep
configuration. This is illustrated in \autoref{fig:fisher_progenitor_overlap_lsst_split}, the LSST
counterpart of \autoref{fig:fisher_progenitor_overlap_split}. In the deep configuration, the
neutron-star fiducial now lies just outside the $1\sigma$ contour of the magnetar model, whereas the
magnetar fiducial still falls inside the $1\sigma$ contour of the neutron-star model, and both remain
within the $2\sigma$ regions. In the shallow configuration, the two models still overlap at the
$1\sigma$ level. A cross-correlation with an LSST-like sample would therefore substantially improve
the prospects for distinguishing between a prompt and a delayed origin of the bursts, even though it would
not yet allow a decisive discrimination.

Beyond a wider footprint, several extensions of the analysis are conceivable.

\begin{itemize}
    \item The parametric FRB redshift distribution of \autoref{eq:frb_nz} could be replaced by
    the measured distribution of an observed catalog; with the second CHIME/FRB
    catalog~\cite{CHIMEFRB2026} a sample of shallow-configuration size is already available.
    \item The Gaussian Fisher formalism could be replaced by a full likelihood sampling, revealing
    non-Gaussian posterior features and allowing a consistent treatment of non-Gaussian covariance
    contributions.
    \item The linear bias model could be replaced by a scale-dependent halo-model description,
    connecting the measured bias to host halo masses and extending the analysis to smaller scales.
\end{itemize}

In conclusion, the results of this thesis demonstrate the strong 
potential of FRB$\times$galaxy cross-correlation as a probe of FRB progenitors. 
Combining the FRB sample with a dense photometric galaxy sample efficiently 
recovers clustering information and enables constraints on the host bias at 
order unity with current survey footprints. With the larger and deeper imaging 
surveys expected in the next generation, this approach could provide sufficient 
precision to distinguish between different progenitor populations and shed light 
on the astrophysical origin of FRBs.


NOTIZEN:

MARKOV CHAIN BETTER FOR PRIORS:
"A full Bayesian analysis with Markov Chain Monte Carlo sampling would be 
required to accurately characterize the posterior in that regime and to assess 
potential biases from prior volume effects."




verlauf bias von frb --> ergaenzen in limitations satz aus slcak von klara und kuerzen ?


Checken der bibliography ob alle links gehen

einmal durchlesen und nochmal alles verstehen

nochmal das correlation chapter angucken und verstehen

Acknowledgements: LsstY10 frau , klara usw..



den 1x2 plot beschreibung checken

bei allen plots die groessen abchecken und ggf. anpassen


Die populaerwissenschaftliche zitat sache --> ausgerechnet quelle oder rechnung in den anhang der BA packen, weil keine quelle gefunden


tabellen angucken und hecken ob werte korerkt sind die dort gekuerzt wurden. eventuell eine nachkomma stelle hinzufuegen bei den sigmas 

TABELLEN NOCHMAL ANGUCKEN UND DIE STRICHEN ANPASSEN, DAMIT SIE BESSER ZUSAMMENPASSEN. AUCH DIE SPALTENBREITEN KÖNNTEN NOCHMAL ANGEPASST WERDEN, DAMIT ES BESSER AUSSIEHT.

ALLE TABELLEN WERTE NOCHMAL CHECKEN

besonders durch anpassung von dem omegab parameter der sich geaendert hat 

checken der quellen in density chapter da durch ai

shot noise der galaxy sample einpflegen.

neutron star merger populatiin einheitlich ueberall benennen

die beschreibung der kreuzspektren nochmal ansehen in method

D. Redundanzen
Introduction und Background wiederholen Lorimer-Burst, CHIME-Zahlen und SGR 1935+2154. Den Lorimer-Absatz in 02background.tex:204 streichen.
„Repeater schließen katastrophale Modelle aus“ steht fast gleich in 02background.tex:227 und 02background.tex:236.
„Optimistic estimates“ steht in 04results.tex:7 und 05conclusion.tex:29.
„CHIME-like/SKA-like“ wird in 04results.tex:224 zum dritten Mal erklärt.
In der Discussion sagt der Satz um L244–246 dreimal dasselbe („comfortably within … well inside … overlap almost entirely“).
Der Jacobian dz/dχ steht zweimal (n(χ) und W(χ)). Die n(χ)-Gleichung kann raus.
In der Conclusion folgt auf „strong constraints“ noch „demonstrate the strong constraining power“.

E. Kürzen
Abschnitte 2.1–2.3 zusammenlegen. Die zweite Friedmann-Gleichung (02background.tex:93) wird nie benutzt. a(t₀) = 1 und H₀ = H(t₀) im Fließtext statt als eigene Gleichungen.
Tabelle tab:fishersetup (04results.tex:52) ist komplett redundant: Die Fiducialwerte stehen schon in tab:frbbias, die Schrittweite von 1 % passt in einen Satz.
KiDS-Details (Teleskop, Seeing, Grenzmagnitude) auf einen Satz reduzieren.

F. Kleinigkeiten (optional)
15th October 2026 → amerikanisch „October 15, 2026“.
In den Plot-Legenden steht „FRB×Galaxy“, im Text „FRB×galaxy“.
